\chapter{Adrenal Fatigue}

\section*{Definition and Overview}
\textit{Adrenal fatigue} is a term used mainly in alternative and complementary medicine to describe a collection of symptoms -- persistent tiredness, difficulty waking in the morning, "brain fog", salt cravings, and a general feeling of being run down -- said to result from adrenal glands that have been "worn out" by chronic stress. It is important to be clear that adrenal fatigue is \textbf{not a recognised medical diagnosis}. Major medical and endocrine bodies have stated that there is no scientific evidence that the adrenal glands can become "fatigued" in this way, and the concept has not been validated by research.

Genuine disorders of the adrenal glands, most notably Addison's disease (primary adrenal insufficiency), are real but quite different conditions that are diagnosed with specific hormone tests and require lifelong treatment. The practical message for anyone troubled by ongoing fatigue is that the correct next step is a proper medical assessment to find the true cause, rather than self-diagnosis of adrenal fatigue or reliance on unvalidated supplements and testing.

\section*{Epidemiology}
Because adrenal fatigue is not a recognised condition, it has no formal epidemiology. The symptom of persistent tiredness itself, however, is extremely common: roughly 1 in 5 adults reports feeling unusually tired, and fatigue is among the most frequent reasons for consulting a general practitioner. In the large majority of people who undergo proper assessment, the fatigue has an identifiable medical, psychological, or lifestyle explanation. Genuine adrenal disease, by contrast, is rare, with primary adrenal insufficiency affecting approximately 1 in 10,000 to 20,000 people, most commonly in adulthood and slightly more often in women.

\section*{Aetiology and Pathophysiology}
The theory behind "adrenal fatigue" proposes that prolonged stress causes the adrenal glands to overwork and then to underproduce hormones such as cortisol. This idea is not supported by evidence. The adrenal glands, which sit above the kidneys, produce cortisol and aldosterone under the control of a finely balanced feedback system involving the pituitary gland and hypothalamus, and they do not simply "run out" under stress.

\begin{itemize}
    \item \textbf{Real causes of fatigue:} sleep disorders such as obstructive sleep apnoea, hypothyroidism, anaemia, diabetes, depression and anxiety, chronic infections, and chronic heart, kidney, or liver disease
    \item \textbf{Lifestyle factors:} chronic sleep deprivation, shift work, excess alcohol, physical inactivity, and obesity
    \item \textbf{Genuine adrenal insufficiency:} a real disease in which the glands fail to produce enough hormones, most often from autoimmune destruction, but also from infections, haemorrhage, or pituitary disease
    \item \textbf{Stress and cortisol:} chronic stress can affect wellbeing and sleep, but it does not "exhaust" the adrenal glands, and cortisol regulation usually remains intact
    \item \textbf{Commercial testing:} saliva or hair cortisol "adrenal fatigue" tests marketed online are not validated for this purpose and frequently produce misleading results
\end{itemize}

\section*{Clinical Features}
The symptoms attributed to adrenal fatigue are non-specific, meaning they overlap substantially with many other common conditions, which is exactly why self-diagnosis is unreliable.

\begin{itemize}
    \item Persistent tiredness and waking unrefreshed despite adequate time in bed
    \item Difficulty getting up in the morning and a pronounced afternoon energy slump
    \item "Brain fog", poor concentration, and memory difficulties
    \item Salt cravings, light-headedness, and feeling worse during stress
    \item Reduced tolerance for stress, irritability, and low mood
    \item By contrast, genuine adrenal insufficiency (Addison's disease) typically causes profound fatigue, weight loss, low blood pressure with dizziness on standing, darkening of the skin, nausea, abdominal pain, and a strong craving for salt
\end{itemize}

\section*{Differential Diagnosis}
Because the presenting symptoms are so non-specific, the real clinical question is what is actually causing the fatigue.

\begin{itemize}
    \item Depression, anxiety, and burnout
    \item Sleep disorders, particularly obstructive sleep apnoea and insomnia
    \item Hypothyroidism and other endocrine disorders
    \item Anaemia from iron, folate, or vitamin B12 deficiency
    \item Diabetes and other metabolic disturbances
    \item Chronic fatigue syndrome (ME/CFS)
    \item Chronic heart, kidney, liver, or lung disease
    \item Medication side effects, including from sedatives and some blood-pressure drugs
    \item Genuine adrenal insufficiency (Addison's disease) -- rare but serious
\end{itemize}

\section*{When to Seek Medical Care}
See a doctor if you have persistent, unexplained tiredness that is affecting daily life, particularly if it has lasted more than a few weeks. Urgent attention is needed if fatigue is accompanied by weight loss, darkening of the skin, dizziness on standing, or a strong craving for salt, because these features may point to genuine adrenal insufficiency.

\textbf{Call 000 (or local emergency services) for:}
\begin{itemize}
    \item Collapse or fainting associated with a very low blood pressure
    \item Severe weakness, vomiting, and confusion in someone with known adrenal disease -- features of an adrenal crisis
    \item Chest pain, severe breathlessness, or any sudden, severe symptom
\end{itemize}

\section*{Diagnosis and Investigations}
A doctor will investigate the real cause of fatigue rather than accept the label of adrenal fatigue.

\begin{itemize}
    \item \textbf{History and examination:} a detailed review of sleep, mood, diet, medicines, alcohol, and lifestyle, together with a general physical examination
    \item \textbf{Blood tests:} full blood count, iron studies, vitamin B12 and folate, thyroid function, fasting glucose, and liver and kidney function
    \item \textbf{Sleep study:} arranged when obstructive sleep apnoea or another sleep disorder is suspected
    \item \textbf{Adrenal testing:} if genuine adrenal insufficiency is suspected, an early-morning blood cortisol level and a short ACTH (Synacthen) stimulation test, arranged through an endocrinologist
    \item \textbf{Not recommended:} "adrenal fatigue" saliva, urine, or hair tests sold commercially, which lack evidence and can cause unnecessary worry or a false sense of security
\end{itemize}

\section*{Management}
The correct management depends entirely on what the assessment actually finds.

\begin{itemize}
    \item \textbf{Treat the identified cause:} thyroid replacement for hypothyroidism, iron supplementation for anaemia, continuous positive airway pressure for sleep apnoea, or treatment of any underlying medical condition
    \item \textbf{Psychological support:} treatment of depression or anxiety with counselling, cognitive behavioural therapy, or medication as appropriate
    \item \textbf{Lifestyle:} regular sleep hours, daily physical activity, a balanced diet, adequate hydration, and reduced alcohol and caffeine intake
    \item \textbf{Stress management:} realistic workloads, relaxation techniques, and maintained social connections
    \item \textbf{Genuine adrenal insufficiency:} lifelong replacement of cortisol and, where needed, aldosterone under specialist care, with written "sick day rules" that advise doubling doses during illness to prevent crisis
    \item \textbf{Healthy scepticism:} be wary of expensive supplements and unvalidated testing promoted for "adrenal fatigue"
\end{itemize}

\section*{Complications}
\begin{itemize}
    \item Missed or delayed diagnosis of a genuine, treatable condition such as hypothyroidism, depression, or sleep apnoea
    \item Delayed recognition of rare but serious adrenal insufficiency, which can progress to a life-threatening adrenal crisis
    \item Unnecessary expense on supplements, consultations, and unvalidated tests
    \item Persistent, untreated symptoms with reduced quality of life and productivity
\end{itemize}

\section*{Prognosis and Outcomes}
When a genuine cause is found and treated, most people improve significantly. Where no specific cause is identified, fatigue often improves over time with attention to sleep, activity, and stress management, although in some people it can be persistent and frustrating. Genuine adrenal insufficiency is a lifelong condition but is very well controlled with replacement therapy, and people with it can lead full and normal lives with careful, specialist-guided management and adherence to sick-day rules.

\section*{Prevention and Self-Care}
\begin{itemize}
    \item \textbf{Sleep:} keep regular sleep hours and protect the quality of your sleep
    \item \textbf{Exercise:} stay active daily, starting gently if you are very tired
    \item \textbf{Diet:} eat a balanced diet and maintain good hydration
    \item \textbf{Stress:} manage workloads realistically and use relaxation and social connection
    \item \textbf{Substances:} limit alcohol and avoid relying on caffeine to push through fatigue
    \item \textbf{Assessment:} see a GP for a proper assessment rather than self-diagnosing
    \item \textbf{Caution:} be critical of health claims for supplements that promise to "support the adrenals"
\end{itemize}

\section*{Sources and Further Reading}
The following sources provide reliable, up-to-date information on fatigue and adrenal disorders.

\begin{itemize}
    \item Mayo Clinic. \textit{Adrenal fatigue: is it real? -- expert commentary on the evidence.} https://www.mayoclinic.org
    \item MSD Manual. \textit{Overview of adrenal insufficiency and Addison disease.} https://www.msdmanuals.com
    \item NHS. \textit{Conditions: tiredness and fatigue.} https://www.nhs.uk/conditions/
    \item MedlinePlus. \textit{Addison disease and adrenal gland disorders.} https://medlineplus.gov
    \item Centers for Disease Control and Prevention. \textit{Health information on fatigue and underlying causes.} https://www.cdc.gov
\end{itemize}
